\documentclass[10pt]{article}
\usepackage[letterpaper,margin=0.64in]{geometry}
\usepackage{booktabs}
\usepackage{graphicx}
\usepackage{array}
\usepackage{tabularx}
\usepackage{xcolor}
\usepackage[hidelinks]{hyperref}
\usepackage{microtype}
\usepackage{enumitem}
\setlist{nosep,leftmargin=1.2em}
\setlength{\parindent}{0pt}
\setlength{\parskip}{2.5pt}
\setlength{\tabcolsep}{3.5pt}
\renewcommand{\arraystretch}{1.06}
\newcommand{\pct}[1]{#1\%}

\begin{document}
\pagestyle{plain}
\begin{center}
{\Large\bfseries Assignment 1: The CNN Challenge}\\[-1pt]
{\normalsize Pritom Kumar Paul \quad---\quad Final Experimental Report}
\end{center}
\vspace{-5pt}

\section*{Final result and validation protocol}
The final model is an ImageNet-pretrained \textbf{ConvNeXt-Large} with all
layers fine-tuned. It achieved \textbf{96.25\% validation accuracy} and
\textbf{96.28\% validation macro F1}. After model selection was locked, a
single labeled-test evaluation produced \textbf{98.50\% accuracy} (394/400)
and \textbf{98.48\% macro F1}.

The 2,400 supplied training images were divided once using seed 0 into a fixed,
class-stratified 1,920/480 train/validation split (120/30 images per class).
SHA-256 grouping kept the six exact duplicate pairs, all in \texttt{LivingRoom},
on the same side of the split. An audit found no exact overlap with either test
directory. The 400 labeled test images were excluded from every architecture,
hyperparameter, and stopping decision; the 400 unlabeled \texttt{test2} images
were unused. Mixed grayscale/RGB source files were all decoded as three-channel
RGB for the final model.

\section*{The final recipe}
Torchvision's ConvNeXt-Large was initialized with \texttt{IMAGENET1K\_V1}
weights; the 1,000-way classifier was replaced by a 16-way head with dropout
0.3. Inputs were 224$\times$224 RGB. Training augmentation used random resized
crop (scale 0.75--1.0, ratio 0.8--1.25), horizontal flip ($p=0.5$), brightness
and contrast jitter (0.15), and random erasing ($p=0.2$). Validation used resize
then center crop. Both used ImageNet mean/std normalization.

All 196.25M parameters were optimized with AdamW (learning rate $10^{-4}$,
weight decay $10^{-4}$), cosine decay, 0.1 label smoothing, batch size 32, AMP,
and a 50-epoch budget with patience-10 validation early stopping. The best
checkpoint occurred at epoch 2; training stopped after epoch 12. On an A100
40GB, the run took 202.3s and peaked at 9,364.6 MiB allocated GPU memory. The
model has an estimated 34.36G MACs (68.72G FLOPs) per image; pooling and
activation operations are excluded from that estimate.

\begin{figure}[h]
\centering
\includegraphics[width=.475\linewidth]{reports/figures/checkpoint6/convnext_large_pretrained_all/loss_curve.png}\hfill
\includegraphics[width=.475\linewidth]{reports/figures/checkpoint6/convnext_large_pretrained_all/accuracy_curve.png}
\vspace{-7pt}
\caption{Checkpoint 6 learning curves. The saved checkpoint is chosen by validation accuracy, not the final epoch.}
\vspace{-5pt}
\end{figure}

\section*{Experimental journey}
\begin{center}
\footnotesize
\begin{tabularx}{\linewidth}{@{}lXrrr@{}}
\toprule
Checkpoint & Controlled question / model & Val. acc. & Params & MACs/img \\
\midrule
1 & Supplied grayscale CNN & 46.04\% & 0.058M & 0.0006G \\
2 & RGB augmentation + regularized scratch ResNet-18 & 76.46\% & 11.18M & 0.59G \\
3 & Scratch family winner: ResNeXt-50 & 78.33\% & 23.01M & 1.38G$^*$ \\
4 & ResNeXt-50, pretrained, last stage + head & 94.79\% & 23.01M & 4.23G \\
5 & ConvNeXt-Tiny, pretrained, all layers & 96.04\% & 27.83M & 4.45G \\
5 & EfficientNet-B0, pretrained, head only & 93.33\% & 4.03M & 0.38G \\
6 & ConvNeXt-Large, pretrained, all layers & \textbf{96.25\%} & 196.25M & 34.36G \\
\bottomrule
\end{tabularx}
\end{center}
\vspace{-3pt}
{\footnotesize $^*$Checkpoint 3 used 128$\times$128 inputs; checkpoints 4--6 used 224$\times$224.}

The starter CNN's widening train/validation gap motivated augmentation and
regularization. Across scratch models, ResNeXt-50 led, but ImageNet pretraining
was the dominant improvement: its four freezing strategies reached
94.17--94.79\%, versus 82.50\% from scratch at 224 pixels. The broader transfer
matrix then selected fully fine-tuned ConvNeXt-Tiny. Scaling to Large gained
only one validation image while using 7.05$\times$ the parameters and
7.71$\times$ the MACs. Large is therefore the final accuracy model, while Tiny
is the stronger accuracy--efficiency choice. EfficientNet-B0 head-only is a
notable low-training-cost alternative: 93.33\% with only 20,496 trainable
parameters.

\section*{Failure analysis and final test behavior}
The clearest failure was scratch MobileNetV3-Large: validation remained at
chance (6.25\%) and every validation image was predicted as \texttt{Suburb},
although training accuracy reached 66.8\%. This exposes a limitation of the
``same recipe for every architecture'' comparison: equal settings are
controlled but not necessarily suitable. The failed run was retained rather
than silently discarded; later comparisons used pretrained representations
and scope-specific learning rates.

On the final test, 12/16 classes had 25/25 correct. \texttt{Kitchen},
\texttt{LivingRoom}, and \texttt{Mountain} each had 24/25; \texttt{TallBuilding}
had 22/25, confused once each with \texttt{Forest}, \texttt{Industrial}, and
\texttt{InsideCity}. The higher test than validation accuracy is plausible for
these small samples and was not used to revise the model.

\section*{Reproducibility, checkpoint, and AI collaboration}
Code and evidence are in the
\href{https://github.com/PritomKumarPaul/computer-vision-assignment1}{GitHub repository},
tag \texttt{checkpoint6}. The exact configuration is
\path{configs/checkpoint6/convnext_large_pretrained_all.yaml}; training is
launched by \path{scripts/run_checkpoint6_convnext_large.sh}. The 749 MiB
checkpoint is the release asset below, followed by its SHA-256:
\begin{center}
{\scriptsize\texttt{checkpoint6\_convnext\_large\_pretrained\_all.pt}}\\[-1pt]
{\tiny\texttt{d011ab54be56d0fbb0c643c1f624d935b1521b1c21f2f73f3df5b2f4649f8919}}
\end{center}
Saved histories, validation matrices, complexity measures, and the immutable
test JSON are tracked in \texttt{results/}.

OpenAI Codex helped inspect the handout/notebook, audit duplicates and color
modes, implement reusable training/evaluation/plotting infrastructure, and
organize evidence. The student ran every experiment, reviewed outputs, chose
the hypothesis sequence and final model, and retained responsibility for the
interpretation. One ineffective AI-assisted choice---using a shared scratch
recipe for MobileNet---is reported above. Full disclosure is in
\texttt{AI\_USAGE.md}.

\end{document}
